\section{Constructive mathematical results}

The original autonomous identity-mass raw-stage target is
\begin{equation}
 (I-h\gamma J_0)U_i=h\gamma f\left(y+\sum_{j<i}a_{ij}U_j\right)
 +\gamma\sum_{j<i}c_{ij}U_j.\label{eq:df-target}
\end{equation}
Both coefficient arrays are strictly lower triangular. Stage increments
have state units, whereas $h$ has time units and $J_0$ inverse-time units.
The following constructions preserve this target or state explicitly the
different fixed-point target they solve.

Parallel iterated Runge--Kutta methods already distinguish the total number
of stages from the sequential iteration depth~\cite{vdH1991}.
Depth-one Anderson extrapolation is a standard construction~\cite{pollock2021}.
The present contribution is the exact reuse/specialization and actual-residual
screen below. For changed frozen operators, order cannot be inherited solely
from the old tableau; this is consistent with the coupled approximation/order
analysis of Rosenbrock--Krylov methods~\cite{tranquilli2014}.

\subsection{Exact reuse of represented causal work}

\paragraph{Theorem A.}
Fix the arithmetic environment, operation order, basis, frozen action,
coefficients, physical state, time, and model epoch. Suppose callbacks are
pure and deterministic. A cache keyed by every bit of
$x_i=\operatorname{fl}(\sum_{j<i}a_{ij}z_j)$ may reuse the represented
projected remainder at stage $i$ without changing any numerical result.
For a fixed block-lower-triangular affine matrix and deterministic diagonal
solve, the maximal unchanged prefix of RHS blocks may likewise reuse its
previous solution prefix.

\paragraph{Proof.}
The complete projected input key fixes the lift, frozen action, physical
state, callback result, subtraction, and projection. For the affine solve,
block zero is identical when its RHS is identical. Induction proves that
an unchanged RHS prefix has an identical represented solution prefix;
the same ordered suffix calculations then reproduce a full solve. This
proves equality of every Picard iterate and residual decision. The full
space check remains unchanged. Budgets may differ because fewer operations
are executed.\hfill$\square$

A physical-state-only key would not necessarily fix the frozen action.
The native cache therefore uses the entire represented projected precursor
and is not carried across a basis or physical-step change.

\subsection{Finite causal propagation and residual transport}

\paragraph{Theorem B.}
For a finite acyclic block map $G_i(x_{<i})$, simultaneous Picard iteration
propagates exact root values by one dependency depth per update. If source
blocks are exact initially, the entire root is exact after at most the
maximum graph depth. A global contraction factor is unnecessary.

\paragraph{Proof.}
The value at a source is independent of the iterate. Once every predecessor
of a node is exact, its next evaluation equals its unique causal-root value.
Induction on graph depth gives the claim.\hfill$\square$

For a chain of $s$ stages, repeated evaluation before each update and at
the final screen costs at most $s(s-1)$ nontrivial callbacks. Exact caching
of settled precursors reduces the corresponding bound to
$s(s-1)/2$. For $s=8$, the independent rational example gives $56$ and $28$.
Native floating behavior is measured using actual identical keys; the
finite exact-arithmetic bound is not substituted for observed work.

\paragraph{Theorem C.}
Suppose the candidate and root are in a domain where
\begin{equation}
 \|G_i(x)-G_i(z)\|\le\sum_{j<i}L_{ij}\|x_j-z_j\|,
 \qquad L_{ij}\ge0.
\end{equation}
With block residual norms $\rho_i=\|G_i(x)-x_i\|$, the block root error obeys
\begin{equation}
 e\le (I-L)^{-1}\rho
   =\sum_{k=0}^{s-1}L^k\rho.\label{eq:df-causal-majorant}
\end{equation}

\paragraph{Proof.}
Adding and subtracting $G_i(x)$ gives $e\le\rho+Le$.
Iterate this inequality $s$ times and use $L^s=0$.\hfill$\square$

For $G_\lambda(x)=(1-\lambda)a+\lambda G(x)$, replace $L$ by
$\lambda L$. Root motion satisfies
\begin{equation}
 |x^*_\mu-x^*_\lambda|_{\rm block}
 \le(\mu-\lambda)(I-\mu L)^{-1}|G(x^*_\lambda)-a|_{\rm block}.
\end{equation}
This provides a finite causal bound even when an unweighted norm of $L$
exceeds one. Whole-domain Lipschitz constants and domain membership are
the required assumptions; a sampled slope does not establish them.

\subsection{Exact two-scalar-tableau physical specialization}

Consider the immutable family
\begin{equation}
 f(y)=b-dy+\kappa P y,\quad Pv=\operatorname{mean}(v)\boldsymbol{1},
 \quad d\ge\kappa\ge0.
\end{equation}
Let $Q=I-P$, $\delta_c=d-\kappa$, and $\delta_t=d$.
The projectors are orthogonal in Euclidean norm and
\begin{equation}
 J=-\delta_cP-\delta_tQ,\qquad
 (I-h\gamma J)^{-1}=
 \frac{P}{1+h\gamma\delta_c}+\frac{Q}{1+h\gamma\delta_t}.
\end{equation}

\paragraph{Theorem D.}
For arbitrary state and forcing vectors, define $f_0=f(y)$ and
\begin{equation}
 t_i(\delta)=\frac{h\gamma+
 \sum_{j<i}(\gamma c_{ij}-h\gamma\delta a_{ij})t_j(\delta)}
 {1+h\gamma\delta}.
\end{equation}
The unique original-stage root is
\begin{equation}
 U_i=t_i(\delta_c)P f_0+t_i(\delta_t)Q f_0.
\end{equation}

\paragraph{Proof.}
Since $f(y+v)=f_0+Jv$, project Eq.~\eqref{eq:df-target} onto the two invariant
subspaces. The projected operator acts by one scalar on each subspace;
causal induction gives the displayed recurrence and lift. Both shifted
denominators are at least one, giving uniqueness.\hfill$\square$

The stage construction costs $O(s^2+sn)$ with no derivative callback,
matrix construction, or Newton step. The retained direct original-space
screen costs an additional $O(s^2n)$ and is charged in both native arms.
For the fixed eight-stage method the complete cost remains linear in $n$.
The original RODAS temporal order is inherited because its exact target and
derivative are preserved. Different floating evaluation orders are compared
with explicit residual and observation tolerances.

For scalar coefficient defects $\epsilon_i(\delta)$ obtained by substituting
the represented coefficients into their recurrence, the exact lifted
stage residual is
\begin{equation}
 \rho_i=\epsilon_i(\delta_c)P f_0+
         \epsilon_i(\delta_t)Q f_0.
\end{equation}
Its squared Euclidean norm separates into the two orthogonal squared terms.
Lifting, initial-RHS, and residual-evaluation rounding require additional
allowances for a rigorous floating enclosure. The current full-space
screen reports an operational residual.

\subsection{Screened vector-only homotopy}

\paragraph{Theorem E.}
Let $G_\lambda:D\to D$ be a contraction of factor $q_\lambda<1$ on a
closed invariant domain. For $r_\lambda(x)=G_\lambda(x)-x$, compute
$p=G_\lambda(x)$. Accept an admissible extrapolate $e$ only when its actual
residual satisfies
\begin{equation}
 \|r_\lambda(e)\|\le\eta\|r_\lambda(p)\|,
 \qquad 0\le\eta\le1.
\end{equation}
Otherwise choose $p$. Then
\begin{equation}
 \|r_\lambda(x^+)\|\le q_\lambda\|r_\lambda(x)\|,
 \qquad \|x-x_\lambda^*\|\le
 \frac{\|r_\lambda(x)\|}{1-q_\lambda}.
\end{equation}

\paragraph{Proof.}
The Picard residual is $G_\lambda(G_\lambda(x))-G_\lambda(x)$, bounded
by $q_\lambda\|r_\lambda(x)\|$. The screen chooses a residual no larger
than that value. Adding and subtracting $G_\lambda(x)$ in the root error
and absorbing the contraction term proves the second claim.\hfill$\square$

The implemented trial is depth-one Anderson extrapolation:
\begin{equation}
 \alpha=\frac{\langle r,r-r_{old}\rangle}{\|r-r_{old}\|^2},\qquad
 e=G_\lambda(x)-\alpha[G_\lambda(x)-G_\lambda(x_{old})].
\end{equation}
Only map evaluations, vectors, and scalar products are used. Degenerate or
inadmissible trials are skipped; rejected evaluated trials are charged.
History is reset at each parameter station. The runtime performs no
Jacobian/JVP/AD request or Newton linear solve.

For the cubic damping implicit-midpoint consumer,
\begin{equation}
 F_\lambda(K)=K+\lambda h[k(y+K/2)+g(y+K/2)^3],
 \quad k>0,\ g\ge0,
\end{equation}
global monotonicity gives $F'_\lambda\ge1+\lambda hk/2$ and hence
\begin{equation}
 |K-K^*_{\lambda}|\le
 \frac{|F_\lambda(K)|}{1+\lambda hk/2}.
\end{equation}
This consumer-specific root bound remains valid when its fixed-point map
is noncontractive. It does not remove midpoint's temporal truncation
error or the need for floating evaluation allowances.

The complete proofs, assumptions, branch-motion and inexact-evaluation
bounds, source comparison, and executed rational/CAS checks are retained
in \texttt{theory/THEORY.md} and its associated receipts.

\subsection{Same-checkpoint entry-pilot proposal reuse}

The adapter's optional repair replaces a replay pilot's cold proposal $h_0$
with
\begin{equation}
 h_{\rm entry}=\min\{H,\max(h_0,h_e)\},
\end{equation}
where $H$ is the same accepted replay span
and $h_e$ is the explicit proposal recorded before advancing that batch.
The native configuration keeps its minimum-step floor and clips the actual
attempt to the endpoint. All ODE data, tolerances, acceptance tests, callback
ceilings, attempt caps, and final same-window cost comparisons are unchanged.

\paragraph{Proposition.}
Under those conditions the pilot retains its original target and
qualification contract. Route selection still requires an actually
completed, qualified, cheaper replay, and all rejected or discarded pilot
work remains charged.

\paragraph{Proof.}
Only a trial-step proposal changes. Each resulting native trial faces the
unchanged error and budget checks; the replay begins at the same state and
ends at the same target. The unchanged route predicate therefore still
requires the original qualification and observed cost inequality. Failed
pilots retain their recorded cost.\hfill$\square$

This preserves the comparison contract, not an adaptive mesh or a promised
speedup. A simple startup bound explains the mechanism: if the first accepted
step is at most $h_s$ and accepted steps grow by at most $g>1$, covering $H$
requires $m$ satisfying
\begin{equation}
 H\le h_s\frac{g^m-1}{g-1}.
\end{equation}
Reusing a larger available proposal reduces this startup lower bound;
actual rejection costs and the accuracy-qualified result are measured.

\subsection{Derived constants for the cubic physical accuracy budget}

Let $\phi(z)=kz+gz^3$, $k>0$, $g\ge0$, and $a=\lambda h/2$.
Writing $m=y+K/2$ gives $F_\lambda(K)=K+2a\phi(m)$ and the unique
root equation $m+a\phi(m)=y$. For $y\ge0$ the midpoint root lies in
$[0,y]$ and the increment root lies in $[-2y,0]$; negative inputs follow
by odd symmetry. On $|m|\le Y$ the sharp derivative bounds are
\begin{equation}
 k\le\phi'(m)\le L_Y=k+3gY^2,
 \qquad F'_\lambda(K)\ge D_\lambda=1+ak.
\end{equation}
The increment map has exact Lipschitz constant $aL_Y$ on that interval.
For $y\ne0$ its analytic root bracket is invariant precisely when
$a(k+gy^2)\le1$; $aL_Y<1$ with $|y|\le Y$ implies both invariance and
contraction. A root bracket alone does not establish a Picard self-map.

Polynomial factorization, for any trial and exact root, gives
\begin{equation}
 F_\lambda(K)=(K-K_*)
 \{1+a[k+g(m^2+mm_*+m_*^2)]\}.
\end{equation}
The quadratic is nonnegative. It proves the global sharp-divisor bound
\begin{equation}
 |K-K_*|\le \frac{|F_\lambda(K)|}{1+ak},
\end{equation}
including trials outside the contraction domain.

\paragraph{Theorem F.}
Let the true solution of $u'=-\phi(u)$ start in $[0,Y]$. Put
\begin{equation}
 \Phi=\phi(Y),\quad L=k+3gY^2,\quad
 C=\frac{L^2\Phi}{8},\quad D=1+hk/2.
\end{equation}
The exact implicit-midpoint map $\Psi_h$ is globally nonexpansive and
has local truncation error at most $\tau=Ch^3/D$.

\paragraph{Proof.}
If $m+(h/2)\phi(m)=y$, then
\begin{equation}
 \Psi_h(y)=2m-y,\qquad
 \Psi'_h(y)=\frac{1-(h/2)\phi'(m)}{1+(h/2)\phi'(m)},
 \qquad |\Psi'_h(y)|\le1.
\end{equation}
For the true increment $\Delta=u(h)-u(0)$, put
$z=[u(0)+u(h)]/2$ and $v=u(h/2)$. Its target residual is
\begin{align}
 F(\Delta)&=A-B,\\
 A&=h[\phi(z)-\phi(v)],\\
 B&=\int_0^h\phi(u(t))\,dt-h\phi(v).
\end{align}
Both defects are nonnegative, since $u$ and $\phi(u)$ are convex.
On the true orbit, $u''\le M_2=L\Phi$ and
$(\phi(u))''\le M_3=6gY\Phi^2+L^2\Phi$. Taylor and midpoint quadrature
remainders give $A\le h^3LM_2/8$ and $B\le h^3M_3/24$. The identity
\begin{equation}
 2L^2-\phi''(Y)\Phi=2k^2+6kgY^2+12g^2Y^4\ge0
\end{equation}
implies $M_3\le3LM_2$, hence $|A-B|\le Ch^3$. The global residual
divisor $D$ proves the local-error bound.\hfill$\square$

The truncation constant $C$ is a proved upper bound, without a claim of
optimality. If a computed endpoint has certified error $a_j$ relative to
$\Psi_h(\widehat y_j)$, nonexpansivity gives
\begin{equation}
 |\widehat y_j-u(jh)|\le j\tau+\sum_{\ell<j}a_\ell.
\end{equation}
For requested physical error $\epsilon$, choose the smallest admitted
dyadic $N$, with exact represented $h=T/N$ and $Nh=T$, such that
\begin{equation}
 E_{\rm tr}=\frac{CT^3}{N(N+kT/2)}<\epsilon,
 \qquad \frac{T[k+3g(Y+\epsilon)^2]}{2N}<1.
\end{equation}
Then assign the entire remaining budget by
\begin{equation}
 a_{\rm step}=\frac{\epsilon-E_{\rm tr}}{N},\qquad
 \rho_{\rm stage}=D a_{\rm step}.
\end{equation}
There is no discretionary fractional budget split. The second condition
is justified inductively by $|\widehat y_j|\le Y+\epsilon$; the true
orbit used in the truncation proof still lies in $[0,Y]$.

The implementation computes these constants outward. For represented
inputs, let $\mathcal R$ enclose the original polynomial residual and
let $0<D_{\rm lo}\le D$. If $z=\operatorname{fl}(\widehat y+K)$ and
$e_{\rm add}$ encloses its addition displacement, then
\begin{equation}
 |z-\Psi_h(\widehat y)|\le
 \frac{\operatorname{mag}(\mathcal R)}{D_{\rm lo}}+e_{\rm add}.
\end{equation}
This bound authorizes an accepted step. Value-only monotone bisection can
polish a candidate; unresolved interval sign or a representability floor
produces a recorded failure. All polishing and rejected trial work is
charged. Observation time rounding is included separately through
$\Phi|\widehat t-jh|$. These are sealed-model certificates, not generic
certificates for an arbitrary callback.

\subsection{Residual dominance with no fitted acceptance constant}

Theorem E holds with $\eta=1$: any admissible candidate whose actual
residual does not exceed the computed Picard candidate's residual
inherits the exact contraction inequality. The new
\texttt{ResidualDominance} mode therefore needs no empirical gain margin,
coefficient cap, observed-rate threshold, or rejection-count threshold.
Normalize the two residual vectors by their common maximum magnitude
before forming the scalar Anderson quotient. A zero represented
denominator, nonfinite coefficient, or inadmissible trial skips the
proposal. No discretionary near-zero epsilon is required; a finite
poorly conditioned proposal must still pass the actual residual screen.
This preserves the selected-update proof but does not guarantee that the
extra map call reduces total work. Floating screening remains operational;
the cubic interval bound above supplies the physical accuracy decision.

\subsection{Sharp accumulated residual gain of the structural quotient}

For a scalar channel rate $r$, define the raw-stage operator
\begin{equation}
 M(r)_{ii}=1+h\gamma r,\qquad
 M(r)_{ij}=h\gamma r a_{ij}-\gamma c_{ij}\quad(j<i).
\end{equation}
Solve $M(r)^T w(r)=b_{\rm code}$ by reverse scalar recurrence. With
$w_t=w(d)$, $w_c=w(d-\kappa)$, define
\begin{equation}
 T_i=w_{t,i}Q+w_{c,i}P,\quad S=R_tQ+R_cP,\quad
 R_r=1-r\sum_i b_i t_i(r).
\end{equation}
For any real scalars $t,c$, put
\begin{equation}
 \nu_n(t,c)=\left|t+\frac{c-t}{n}\right|
 +(n-1)\left|\frac{c-t}{n}\right|.
\end{equation}

\paragraph{Theorem G.}
The sharp one-step stage-residual and propagation infinity norms are
\begin{equation}
 A_n=\sum_i\nu_n(w_{t,i},w_{c,i}),\qquad
 G_n=\nu_n(R_t,R_c).
\end{equation}
For $N$ equal steps of the same autonomous model, the sharp accumulated
norm from independent original-stage residuals is
\begin{equation}
 C_N=\sum_{j=0}^{N-1}\sum_i
       \nu_n(R_t^j w_{t,i},R_c^j w_{c,i}).
\end{equation}
If all those residuals have infinity norm at most $\delta$ and the
initial states agree, the algebraic discrete-trajectory error at every
grid index $m\le N$ is at most $C_N\delta$. Thus, for positive gain,
$\delta=\epsilon_{\rm alg}/C_N$ derives the residual tolerance.

\paragraph{Proof.}
The exact stage-error equation is $M(r)E=R$, so the endpoint error is
$b_{\rm code}^TE=w(r)^TR$ in each channel. The homogeneous propagation
is $S$ by the exact scalar quotient. A matrix $tQ+cP$ has diagonal
$t+(c-t)/n$ and off-diagonal $(c-t)/n$, giving the absolute row sum
$\nu_n$. With zero initial error,
\begin{equation}
 e_N=\sum_{\ell=0}^{N-1}\sum_i
       S^{N-1-\ell}T_iR_{\ell,i}.
\end{equation}
The projector identities give
$S^jT_i=R_t^jw_{t,i}Q+R_c^jw_{c,i}P$. Summing the absolute row
entries of the concatenated independent inputs gives $C_N$. Choosing
each input sign to align with a fixed output row attains equality;
invertibility of the triangular stage systems realizes those residuals.
Nonnegativity of the terms gives $C_m\le C_N$.\hfill$\square$

Submultiplicativity also yields
$C_N\le A_n\sum_{j=0}^{N-1}G_n^j$, without assuming $G_n<1$.
Equality can hold, including in the present audited fixture. Scalar
powers and finite sums avoid a singular expression at $G_n=1$.
The native budget evaluates the formula in binary64 and reports an
operational derived tolerance. It does not include temporal truncation,
endpoint rounding, coefficient error, or an interval residual enclosure.

\subsection{Exact integer preference for enabled entry pilots}

Use the existing work objective before floating conversion,
\begin{equation}
 W=N_{\rm RHS}+N_{\rm JVP}+N_{f_t}+nN_J+N_{\rm solve}.
\end{equation}
The optional strict policy accepts a successful same-checkpoint,
same-endpoint pilot precisely when its checked integer $W$ is smaller
than the accepted explicit prefix's $E$. Ties retain the explicit route.
Any running callback-prefix count $L\ge E$ rules out a strict win,
because final work is at least $L$. Since the native guard charges the
crossing callback and stops on $L>\mathrm{cap}$, the exact early ceiling
is $E-1$ for $E>0$. Integer overflow or an unrepresentable ceiling
declines that evidence instead of truncating the count. All attempted
pilot work remains charged.

This replaces a fitted entry margin only for enabled pilots. Zero pilot
budget retains the existing pressure-only direct-entry mode, and the
pressure, retry, and return heuristics remain separate. The conclusion
is a local choice in the stated work objective; whole-trajectory and
wall-clock gains are measured separately.
